\chapter{Introduction}
\label{ch:introduction}

\section{Motivation}
\label{sec:intro-motivation}

Machine learning and quantum computing are, individually, two of the most
consequential computational developments of the past two decades. Their
intersection --- \emph{quantum machine learning} (QML) --- has attracted
extraordinary attention, driven by the prospect that quantum computers might
learn from data in ways that classical computers cannot
\citep{biamonte2017qml, schuld2021book}. Parameterized quantum circuits can
represent functions in Hilbert spaces whose dimension grows exponentially with
the number of qubits, and quantum kernels can measure similarities between data
points embedded in those spaces \citep{havlicek2019kernels, schuld2019hilbert}.
On paper, this looks like a recipe for more expressive models.

The empirical record is far more ambiguous. The devices available today ---
noisy intermediate-scale quantum (NISQ) hardware \citep{preskill2018nisq} ---
support only small, shallow circuits. Provable quantum--classical learning
separations exist but are confined to carefully constructed problems
\citep{liu2021rigorous}, and a growing body of work shows that classical models
given access to the same data often match or exceed quantum models on the
benchmarks actually used in the literature \citep{huang2021power,
bowles2024better, schnabel2025scrutiny}. At the same time, methodological
audits of published QML studies find recurring flaws that systematically favor
the quantum contender: tuned quantum models compared against default-configured
classical baselines, cherry-picked datasets, inconsistent data splits, and
metrics aggregated across incompatible execution regimes
\citep{bowles2024better}.

This thesis takes the tension between promise and evidence as its starting
point. Rather than asking whether quantum machine learning is ``better'' in the
abstract, it asks a narrower, answerable question: \emph{under carefully matched
conditions --- identical data splits, identical preprocessing, identical
hyperparameter-search budgets, and parameter-matched architectures --- how do
quantum and hybrid models actually compare to strong classical baselines on
standard benchmark tasks, and which design factors drive the differences?}
Answering this requires as much care in experimental methodology as in quantum
theory, and the methodological framework is itself one of the main
contributions of this work.

\section{Research questions}
\label{sec:intro-rqs}

Five research questions drive both the theoretical treatment and the
experimental programme of this thesis.

\begin{description}
  \item[RQ1 (Performance).] Under matched data, preprocessing, and
  hyperparameter-search budgets, how do quantum models --- quantum kernels and
  variational quantum classifiers --- compare to strong classical baselines in
  accuracy, balanced accuracy, $F_1$, and ROC-AUC?

  \item[RQ2 (Design factors).] How do the data-encoding strategy, ansatz
  choice, circuit depth, and entanglement structure affect quantum-model
  performance and trainability?

  \item[RQ3 (Hybrid value).] Do hybrid architectures --- classical feature
  extraction feeding a quantum head, quantum kernels inside a classical SVM,
  quantum transfer learning --- outperform either pure approach at matched
  parameter counts?

  \item[RQ4 (Robustness and cost).] How do finite sampling (shot noise),
  simulated device noise, and real quantum hardware affect the results, and
  what are the actual resource costs --- wall-clock time, parameter counts,
  circuit depth, and shot budgets?

  \item[RQ5 (Data regime).] Does the relative performance of classical,
  quantum, and hybrid models change with the training-set size and the feature
  dimensionality?
\end{description}

\section{Scope and honest framing}
\label{sec:intro-framing}

A thesis-scale experimental programme can reach system sizes of roughly ten to
twenty qubits. Every quantum model studied here can therefore be simulated
classically --- indeed, most experiments in this thesis \emph{are} exact
statevector simulations, by design. It follows that \textbf{this thesis cannot,
even in principle, demonstrate quantum advantage}: any task its quantum models
solve is, by construction, solvable classically at the same scale. What it can
do --- and what the strongest recent benchmark literature argues is currently
the most valuable contribution \citep{bowles2024better, schnabel2025scrutiny}
--- is measure the \emph{empirical competitiveness and behavior of quantum
models at NISQ scale} under a protocol fair enough that the observed
differences mean something. Claims in this thesis are therefore claims about
performance, trainability, robustness, and cost at the scales studied, never
about asymptotic separation. Where the literature offers provable separations
\citep{liu2021rigorous} or arguments about the role of data access
\citep{huang2021power}, they are treated in Chapter~\ref{ch:related} as context
for interpreting the empirical results, not as claims this thesis tests at
scale.

The scope is further restricted to \emph{supervised binary classification} on
standard benchmark datasets --- synthetic constructions with controlled
structure, small tabular datasets, and reduced image data --- the setting in
which the comparative literature is concentrated and in which matched
conditions are achievable within a Master's project. Multi-class extensions,
regression, generative models, and learning on inherently quantum data are
discussed as outlook.

\section{Contributions}
\label{sec:intro-contributions}

The main contributions of this thesis are:

\begin{enumerate}
  \item \textbf{A fairness-first experimental protocol}, specified before any
  quantum model was trained: all models --- classical, quantum, and hybrid ---
  share identical stored cross-validation splits (5-fold stratified, five
  seeds), identical $[0,\pi]$ feature scaling fit on training folds only,
  an identical hyperparameter-optimization budget (one 50-trial Optuna study
  per model--dataset pair under a fixed procedure), parameter-matched
  neural-network baselines, and strict separation of execution regimes (exact
  simulation, finite shots, noisy simulation, hardware) in all reported
  results.

  \item \textbf{A systematic empirical comparison} of seven classical
  baselines, quantum kernel methods, variational quantum classifiers, and
  hybrid architectures across ten benchmark datasets spanning saturated to
  genuinely hard, including controlled synthetic tasks (interleaved half-moons,
  concentric circles, and a sign-parity task with distractor dimensions
  designed to separate model families).

  \item \textbf{Focused design-factor studies}: a data-encoding ablation
  connected to the Fourier-expressivity theory of parameterized circuits
  \citep{schuld2021fourier}; a trainability study reproducing barren-plateau
  scaling \citep{mcclean2018barren}; sample-efficiency learning curves; and a
  noise ladder from exact simulation through calibrated device-noise models to
  validation runs on IBM quantum hardware.

  \item \textbf{A fully reproducible, open research artifact}: every figure and
  table in this document is generated by a version-controlled script from
  append-only result logs carrying configuration hashes, seeds, and split
  identifiers; the stored splits and locked hyperparameter configurations are
  part of the repository.
\end{enumerate}

\section{Thesis outline}
\label{sec:intro-outline}

Chapter~\ref{ch:classical} reviews the classical machine-learning foundations
the comparison rests on, with emphasis on kernel methods and evaluation
methodology. Chapter~\ref{ch:quantum} introduces the quantum-computing
concepts required for the remainder: states, gates, measurement, and the
realities of NISQ devices. Chapter~\ref{ch:qml} develops quantum machine
learning in depth --- data encodings and their expressivity consequences,
variational quantum algorithms, quantum kernel methods, the main circuit
architectures, trainability obstacles, and hybrid designs.
Chapter~\ref{ch:related} surveys comparative studies and the quantum-advantage
debate, positioning this thesis within that literature.
Chapter~\ref{ch:methodology} specifies the datasets, the model zoo, and the
fairness protocol in full detail. Chapter~\ref{ch:experiments} reports the
experimental results, one experiment family at a time, each stated as a
hypothesis before its outcome. Chapter~\ref{ch:discussion} synthesizes the
findings across research questions and examines threats to validity, and
Chapter~\ref{ch:conclusion} concludes with directions for future work.
